\section{The degree-five example and a single inertia group}
\label{sec:genus-two-coefficient-two}
The degree-five case exhibits the integral normalization in the smallest
higher-genus curve. Put \(O=\mathbf Z_2[\zeta_5]\), let
\(\sigma(\zeta_5)=\zeta_5^2\), and write \(\bar z=\sigma^2z\).
The CM type is \(\{1,\sigma\}\), its inverse is
\(1+\sigma^{-1}\), and the integral factor is \(Q_5=-\sigma\).
Thus five satisfies the unit criterion. Set
\[
 T_5^{(2)}=\{z\in1+2O:z\bar z=1\},\qquad
 P_5=(\sigma-1)T_5^{(2)},\qquad
 R_{{\rm quad},5}=R_{\rm quad}\cap\ker\mathfrak h_5,
\]
\[
 V_{{\rm quad},5}=V_{\rm quad}\cap\ker\mathfrak h_5.
\]
There is only one prime above two, so
Proposition~\ref{prop:cm-commutator-torus-torsion} gives
\(P_5\simeq\mathbf Z_2^2\).

\begin{corollary}\label{thm:genus-two-exact-image}
The degree-five homology scalar has exact image \(P_5\) on
\(R_{\rm quad}\), \(V_{\rm quad}\), and both arithmetic intersections.
The four corresponding quotients by its kernel are \(\mathbf Z_2^2\).
Moreover \(G=R_{{\rm quad},5}H\); the full relative quotient on
\(R_{{\rm quad},5}\) is \(B\), and its maximal pro-two version on
\(V_{{\rm quad},5}\) is \(B_{\max2}\). Both original restriction
kernels are exactly \(\mathcal S\), and both images are the full
relative invariant character space \(\operatorname{Hom}_{\rm cts}(B,\F_2)\).
\end{corollary}
\begin{proof}
Apply Theorem~\ref{thm:simultaneous-cyclotomic-cm} with \(\Pi=\{5\}\).
All its hypotheses were verified explicitly above. The torus has free
rank two and no sign-torsion image, giving the stated quotient.
\end{proof}

The arithmetic realization can be made inside one inertia group before
taking the final commutator.
\begin{lemma}\label{lem:genus-two-inertia-M}
Every \(v\in T_5^{(2)}\) is \(\mathfrak h_5(j)\) for a pro-two
arithmetic element \(j\in I_2\cap M\).
\end{lemma}
\begin{proof}
Let \(J\) be the Jacobian of \(t^2-t=w^5\). By
Lemma~\ref{lem:pdl-reflex-and-local}, the exact local CM formula is
\(r(u)=u\sigma^{-1}(u)\) on \(O^\times\). To prescribe \(v\),
choose a unit \(z\) with \(\bar z/z=v\), by unit Hilbert 90.
Set \(u=\sigma^{-1}(z)/z\). Then
\[
 N_{O/\mathbf Z_2}(u)=1,\qquad
 u\sigma^{-1}(u)=\bar z/z=v.
\]
An inertia element \(b\) representing this local unit therefore has
Jacobian value \(v\). The Weil pairing gives
\(\lambda_2(b)=v\bar v=1\); all odd cyclotomic components are one
on inertia at two. Thus \(b\in M\). The based homology value may
have an additional deck factor in \(\mu_5\). Taking
\(j=b^{e_2}\), with the two-primary idempotent exponent, kills this
factor and leaves \(v\) unchanged. It preserves inertia and \(M\),
which proves the assertion.
\end{proof}

For inertia \(a\) at five with deck exponent two, the commutator
\([a,j]\) has scalar \(\sigma(v)/v\). Every Kummer coordinate is
zero because \(j\) is pro-two and \(\lambda_2(a)=1\).
Its primary word actions are trivial: the coefficient-two action kills
\(a\), and every odd-primary IA action kills the pro-two element
\(j\). Both entries fix the two elliptic CM fields, so their corrected
CM commutators are one; the cubic Kummer correction is zero on the
rational-radical kernel. Finally all chosen cup characters extend
invariantly to \(M\), and therefore kill \([a,j]\). This gives a
single-commutator realization in \(V_{\rm quad}\cap H\).
The degree-seven example in Section~\ref{sec:prime-degree-local-cm}
gives a further quotient \(\mathbf Z_2^3\) after imposing this kernel.
